%!TEX root = ../main.tex
\chapter{Order-of-Magnitude Estimation}
\label{ch:estimation}

Physicist Enrico Fermi was famous for estimating quantities he had no
direct data for — the number of piano tuners in Chicago, say — by
breaking the question into smaller pieces that could each be guessed
within a factor of two or three, and multiplying them together. The
errors in the pieces tend to partially cancel, and the final estimate
usually lands within an order of magnitude of the truth.

\section{A Worked Example}
\label{sec:worked-example}

\begin{example}[How many leaves are on a large oak tree?]
Estimate: a mature oak canopy is roughly a sphere 15 metres across.
A branch tip carries perhaps 20 leaves in a cluster the size of a
fist, and clusters are spaced roughly 10cm apart across the outer
surface of the canopy. The canopy's surface area is about
$4\pi (7.5\,\text{m})^2 \approx 700\,\text{m}^2$. At one cluster per
$0.01\,\text{m}^2$ of surface, that's about 70{,}000 clusters, and
$70{,}000 \times 20 \approx 1.4$ million leaves — comfortably in the
range (100{,}000 to a few million) that direct counts of oak trees
actually report.
\end{example}

\section{Why the Errors Cancel}
\label{sec:why-errors-cancel}

Each individual guess in the chain above could easily be off by a
factor of two. But an overestimate in one factor is roughly as likely
as an underestimate, and independent errors of that size tend not to
all point the same direction at once. This is why a chain of five
rough guesses, each uncertain by a factor of two, typically lands the
final answer within a factor of three or four of the truth — far
better than the naive intuition that the errors should simply
multiply together.
